%% Generated by docs/build_study_variants.py from backend/seed_study.py and
%% backend/code_mixed.py. Edit those, then regenerate — do not edit by hand.
%% Compile with XeLaTeX (twice): xelatex study_variants.tex
\documentclass[11pt,a4paper]{article}

\usepackage[margin=2.3cm]{geometry}
\usepackage{fontspec}
\usepackage{xcolor}
\usepackage{listings}
\usepackage{booktabs}
\usepackage{tabularx}
\usepackage{longtable}
\usepackage{array}
\usepackage{enumitem}
\usepackage{textcomp}
\usepackage{float}
\usepackage[hidelinks]{hyperref}

\setmainfont{Latin Modern Roman}
\setmonofont{Latin Modern Mono}[Scale=0.95]
\newfontfamily\tefont{Nirmala UI}[Script=Telugu]
\newfontfamily\tafont{Nirmala UI}[Script=Tamil]
\newfontfamily\hifont{Nirmala UI}[Script=Devanagari]
\newfontfamily\gufont{Nirmala UI}[Script=Gujarati]

\setlength{\emergencystretch}{3em}
\setlength{\parindent}{0pt}
\setlength{\parskip}{4pt}

\definecolor{codebg}{RGB}{247,247,247}
\definecolor{codeframe}{RGB}{210,210,210}
\definecolor{kw}{RGB}{0,70,140}
\lstset{
  language=Python,
  basicstyle=\ttfamily\small,
  keywordstyle=\color{kw}\bfseries,
  numbers=left,
  numberstyle=\tiny\color{gray},
  numbersep=6pt,
  backgroundcolor=\color{codebg},
  frame=single,
  rulecolor=\color{codeframe},
  xleftmargin=1.5em,
  framexleftmargin=1.2em,
  showstringspaces=false,
  columns=fullflexible,
  keepspaces=true,
  breaklines=true,
  aboveskip=6pt,
  belowskip=6pt,
}

\title{Same Code, Different Contexts:\\The Eight Study Variants}
\author{COMPUTE 2026 Working Group on Refute Problem Difficulty}
\date{\today}

\begin{document}
\maketitle

\begin{abstract}
This document lists the eight Refute Problem variants used in the study: two
problems, each written in a direct and a contextualized form and paired with
a context-independent and a context-sensitive bug. For every variant it gives
the specification the student reads, the program they are asked to refute,
why the form of the specification may matter for that bug, a checked
counter-example, and the rule a Step~3 input must satisfy. It closes by
explaining how a problem is stored in the study application's database. All
programs, inputs and outputs are taken from, and computed by, the
application itself.
\end{abstract}

\tableofcontents
\newpage

\section{Design in brief}

Each problem has one reference solution and is written in two forms. The
\emph{direct} form states the computation in computational terms; the
\emph{contextualized} form places the same computation in a real-world
setting. Both forms are fully specified and are solved by the same reference
solution. Each problem is paired with two buggy programs: a
\emph{context-independent} bug (Buggy1), whose failing inputs are no easier to
think of in the setting, and a \emph{context-sensitive} bug (Buggy2), which the
setting may help students find, either by suggesting realistic inputs that
expose it or by giving the faulty code a meaning that looks wrong.
Crossing form with bug gives the eight variants of Table~\ref{tab:variants}.

Each student attempts both problems, one in each form, and sees one of each
problem's two buggy programs. Which form, which program, and which problem
comes first are counterbalanced (Table~\ref{tab:conditions}): any four
consecutive rows show every form--bug combination of each problem once, and
all eight also balance the order.

Within a problem the student works through three steps, with no feedback
between them:
\begin{description}[style=nextline, leftmargin=1.5em]
  \item[Step 1: the task.] The student reads the specification (switching to a
    code-mixed version if they wish) and gives three input--output pairs.
    No code is shown. The pairs must all be correct to continue.
  \item[Step 2: the code.] The student sees the buggy program, without the
    specification, and traces it on inputs we supply. None of these inputs is
    a counter-example.
  \item[Step 3: refutation.] The student sees the specification and the program
    together, with their own Step~1 pairs and Step~2 traces, and gives an input
    on which the program fails, the output the task requires for it, and the
    output the program returns. They have at most three attempts.
\end{description}

\subsection{Participant questionnaire}\label{sec:questionnaire}
Before the first problem, every student answers the following on the start
page. All items are required; the stored value is given in brackets where it
differs from what the student sees.
\begin{enumerate}[leftmargin=2em, itemsep=1pt]
  \item Institute (chosen from a list) and participant / student ID.
  \item Your native language: Telugu, Tamil, Hindi, Gujarati or German
    [\texttt{te}, \texttt{ta}, \texttt{hi}, \texttt{gu}, \texttt{de}]. This
    also chooses the code-mixed version of the specification.
  \item Medium of instruction at school: English / My native language /
    Partly English, partly my native language / Other.
  \item Proficiency in English (1 = novice, 5 = expert).
  \item Proficiency in your mother language (1 = novice, 5 = expert).
  \item Proficiency in computer programming (1 = novice, 5 = expert).
  \item How many computer programming courses have you taken? (a whole
    number, 0--50)
  \item Age (years, 10--100).
  \item Gender: Female / Male / Other / Prefer not to say
    [\texttt{NoResponse}].
\end{enumerate}
A student who returns with the same institute and ID resumes their earlier
session; their first answers are kept.
\begin{table}[H]\centering\small
\caption{The eight variants, with one counter-example for each bug.}
\label{tab:variants}
\begin{tabularx}{\textwidth}{@{}l l l X@{}}\toprule
\textbf{Variant} & \textbf{Form} & \textbf{Bug} & \textbf{Counter-example} \\\midrule
\texttt{P1\_Direct\_buggy1} & Direct & Buggy1 (context-independent) & \texttt{has\_large\_change([70, 72, 110], 15)} \\
\texttt{P1\_Contextualized\_buggy1} & Contextualized & Buggy1 (context-independent) & \texttt{has\_large\_change([70, 72, 110], 15)} \\
\texttt{P1\_Direct\_buggy2} & Direct & Buggy2 (context-sensitive) & \texttt{has\_large\_change([72, 74], 15)} \\
\texttt{P1\_Contextualized\_buggy2} & Contextualized & Buggy2 (context-sensitive) & \texttt{has\_large\_change([72, 74], 15)} \\
\texttt{P2\_Direct\_buggy1} & Direct & Buggy1 (context-independent) & \texttt{warmest([-9.8, -12.0, -9.8])} \\
\texttt{P2\_Contextualized\_buggy1} & Contextualized & Buggy1 (context-independent) & \texttt{warmest([-9.8, -12.0, -9.8])} \\
\texttt{P2\_Direct\_buggy2} & Direct & Buggy2 (context-sensitive) & \texttt{warmest([-14.2, -9.8])} \\
\texttt{P2\_Contextualized\_buggy2} & Contextualized & Buggy2 (context-sensitive) & \texttt{warmest([-14.2, -9.8])} \\
\bottomrule\end{tabularx}\end{table}
\begin{table}[H]\centering\small
\caption{Condition rows. Participants at each institute are dealt these rows in shuffled blocks of eight.}
\label{tab:conditions}
\begin{tabular}{@{}r l l l l l@{}}\toprule
\textbf{Row} & \textbf{P1 form} & \textbf{P1 bug} & \textbf{P2 form} & \textbf{P2 bug} & \textbf{Order} \\\midrule
0 & Direct & Buggy1 & Contextualized & Buggy1 & P1 then P2 \\
1 & Direct & Buggy2 & Contextualized & Buggy2 & P1 then P2 \\
2 & Contextualized & Buggy1 & Direct & Buggy2 & P1 then P2 \\
3 & Contextualized & Buggy2 & Direct & Buggy1 & P1 then P2 \\
4 & Direct & Buggy1 & Contextualized & Buggy1 & P2 then P1 \\
5 & Direct & Buggy2 & Contextualized & Buggy2 & P2 then P1 \\
6 & Contextualized & Buggy1 & Direct & Buggy2 & P2 then P1 \\
7 & Contextualized & Buggy2 & Direct & Buggy1 & P2 then P1 \\
\bottomrule\end{tabular}\end{table}
\clearpage
\section{Problem P1: Threshold (heart-rate monitor)}\label{sec:p1}
\subsection{Shared by all four variants}
The function signature, the reference solution, and the inputs used in Steps~1 and~2 are identical in all four variants of this problem. The Step~2 inputs are the same for both buggy programs, and neither program fails on any of them.

\paragraph{Signature.} \texttt{def has\_large\_change(readings, T):}

\begin{lstlisting}[title={\small Reference solution (never shown to students)}]
def has_large_change(readings, T):
    for i in range(len(readings) - 1):
        if abs(readings[i+1] - readings[i]) > T:
            return True
    return False
\end{lstlisting}
\paragraph{Step 2 tracing inputs.} The control-flow question (how many times each line runs) uses the first input; the data-flow question (variable values) uses the second. Reference outputs:

\smallskip\noindent\begin{tabular}{@{}l l@{}}\toprule
\textbf{Input} & \textbf{Output} \\\midrule
\texttt{has\_large\_change([3, 5, 4], 5)} & \texttt{False} \\
\texttt{has\_large\_change([2, 10, 11], 5)} & \texttt{True} \\
\bottomrule\end{tabular}

\paragraph{Step 1 worked examples.} Revealed, with their correct outputs, only to a student who asks for help after a wrong Step~1 attempt; the student may not then reuse them as their own pairs:

\smallskip\noindent\begin{tabular}{@{}l l@{}}\toprule
\textbf{Input} & \textbf{Output} \\\midrule
\texttt{has\_large\_change([70, 95, 96], 15)} & \texttt{True} \\
\texttt{has\_large\_change([], 10)} & \texttt{False} \\
\bottomrule\end{tabular}

\subsection{Variant \texttt{P1\_Direct\_buggy1}}
\textbf{Form:} Direct. \quad \textbf{Bug:} Buggy1, context-independent.

\paragraph{Specification shown to the student (49 words).}
\begin{quote}\small Write a function \texttt{has\_large\_change(readings, T)} that takes a list of integer \texttt{readings} and a positive integer \texttt{T}. It returns \texttt{True} if the absolute difference between some pair of consecutive elements of \texttt{readings} is greater than \texttt{T}, and \texttt{False} otherwise. The list may be empty or contain a single element.\end{quote}
\begin{lstlisting}[title={\small Program shown in Steps 2 and 3}]
def has_large_change(readings, T):
    for i in range(len(readings) - 1):
        if abs(readings[i+1] - readings[i]) > T:
            return True
        else:                                      # <- bug
            return False                           # <- bug
    return False
\end{lstlisting}
\paragraph{The bug.} The \texttt{else} branch returns \texttt{False} as soon as the first pair is within \texttt{T}, so the loop never looks past the first pair. A counter-example needs an unremarkable first pair followed by a large change.

\paragraph{Why the form matters.} The requirement the bug violates (\emph{some} pair, \emph{ever} changes) is stated in this form exactly as in the contextualized one.

\paragraph{Counter-example.} \texttt{has\_large\_change([70, 72, 110], 15)}: the task requires \texttt{True}, the program returns \texttt{False}.

\paragraph{Valid Step 3 inputs.} Inputs must satisfy the specification's own preconditions.
\begin{lstlisting}[title={\small Validity check (server-side)}]
def valid(readings, T):
    if not isinstance(readings, list) or not all(type(r) is int for r in readings):
        return 'readings must be a list of whole numbers.'
    if type(T) is not int or T <= 0:
        return 'T must be a positive whole number.'
    return None
\end{lstlisting}
\paragraph{Expected effect.} Baseline for this bug.

\subsection{Variant \texttt{P1\_Contextualized\_buggy1}}
\textbf{Form:} Contextualized. \quad \textbf{Bug:} Buggy1, context-independent.

\paragraph{Specification shown to the student (127 words).}
\begin{quote}\small A hospital's patient monitor records a patient's heart rate, in beats per minute (bpm), once every second. A healthy heart rate changes only slightly from one second to the next. The hospital's clinical guidelines say that a nurse must review a recording if the heart rate ever changes by more than \texttt{T} bpm between two consecutive readings: a sudden drop can mean the heart is failing, and a sudden rise can mean the patient is in distress. A recording may be empty or contain a single reading, for example if the monitor was just attached. Write a function \texttt{has\_large\_change(readings, T)} that takes the list of readings (whole numbers) and the threshold \texttt{T} (a positive whole number), and returns \texttt{True} if the recording needs review, and \texttt{False} otherwise.\end{quote}
\begin{lstlisting}[title={\small Program shown in Steps 2 and 3}]
def has_large_change(readings, T):
    for i in range(len(readings) - 1):
        if abs(readings[i+1] - readings[i]) > T:
            return True
        else:                                      # <- bug
            return False                           # <- bug
    return False
\end{lstlisting}
\paragraph{The bug.} The \texttt{else} branch returns \texttt{False} as soon as the first pair is within \texttt{T}, so the loop never looks past the first pair. A counter-example needs an unremarkable first pair followed by a large change.

\paragraph{Why the form matters.} Nothing about heart rates makes a failing input more or less likely to occur to a student, and the violated requirement is stated as in the direct form. The setting only adds reading load.

\paragraph{Counter-example.} \texttt{has\_large\_change([70, 72, 110], 15)}: the task requires \texttt{True}, the program returns \texttt{False}.

\paragraph{Valid Step 3 inputs.} Inputs must also be \emph{plausible in the setting}; students are told so, and an implausible Step~3 input is rejected with the message below and uses up an attempt.
\begin{lstlisting}[title={\small Validity check (server-side)}]
def valid(readings, T):
    if not isinstance(readings, list) or not all(type(r) is int for r in readings):
        return 'readings must be a list of whole numbers.'
    if type(T) is not int or T <= 0:
        return 'T must be a positive whole number.'
    if any(r < 30 or r > 220 for r in readings) or T > 60:
        return ('Your input must be a realistic recording: every heart rate '
                'between 30 and 220 bpm, and T at most 60 bpm.')
    return None
\end{lstlisting}
\paragraph{Expected effect.} No gain over Direct; possibly lower for less fluent readers. Longer time (more to read).

\subsection{Variant \texttt{P1\_Direct\_buggy2}}
\textbf{Form:} Direct. \quad \textbf{Bug:} Buggy2, context-sensitive.

\paragraph{Specification shown to the student (49 words).}
\begin{quote}\small Write a function \texttt{has\_large\_change(readings, T)} that takes a list of integer \texttt{readings} and a positive integer \texttt{T}. It returns \texttt{True} if the absolute difference between some pair of consecutive elements of \texttt{readings} is greater than \texttt{T}, and \texttt{False} otherwise. The list may be empty or contain a single element.\end{quote}
\begin{lstlisting}[title={\small Program shown in Steps 2 and 3}]
def has_large_change(readings, T):
    prev = 0                         # <- bug
    for r in readings:
        if abs(r - prev) > T:
            return True
        prev = r
    return False
\end{lstlisting}
\paragraph{The bug.} \texttt{prev = 0} compares the first reading with a value that was never recorded. The program fails only when the first reading itself exceeds \texttt{T} while no real consecutive pair does.

\paragraph{Why the form matters.} Here 0 looks like a natural starting value, and the program is correct on the small inputs a student is likely to try: on \texttt{[1, 2, 3]} with $T=5$ it correctly returns \texttt{False}.

\paragraph{Counter-example.} \texttt{has\_large\_change([72, 74], 15)}: the task requires \texttt{False}, the program returns \texttt{True}.

\paragraph{Valid Step 3 inputs.} Inputs must satisfy the specification's own preconditions.
\begin{lstlisting}[title={\small Validity check (server-side)}]
def valid(readings, T):
    if not isinstance(readings, list) or not all(type(r) is int for r in readings):
        return 'readings must be a list of whole numbers.'
    if type(T) is not int or T <= 0:
        return 'T must be a positive whole number.'
    return None
\end{lstlisting}
\paragraph{Expected effect.} Baseline for this bug.

\subsection{Variant \texttt{P1\_Contextualized\_buggy2}}
\textbf{Form:} Contextualized. \quad \textbf{Bug:} Buggy2, context-sensitive.

\paragraph{Specification shown to the student (127 words).}
\begin{quote}\small A hospital's patient monitor records a patient's heart rate, in beats per minute (bpm), once every second. A healthy heart rate changes only slightly from one second to the next. The hospital's clinical guidelines say that a nurse must review a recording if the heart rate ever changes by more than \texttt{T} bpm between two consecutive readings: a sudden drop can mean the heart is failing, and a sudden rise can mean the patient is in distress. A recording may be empty or contain a single reading, for example if the monitor was just attached. Write a function \texttt{has\_large\_change(readings, T)} that takes the list of readings (whole numbers) and the threshold \texttt{T} (a positive whole number), and returns \texttt{True} if the recording needs review, and \texttt{False} otherwise.\end{quote}
\begin{lstlisting}[title={\small Program shown in Steps 2 and 3}]
def has_large_change(readings, T):
    prev = 0                         # <- bug
    for r in readings:
        if abs(r - prev) > T:
            return True
        prev = r
    return False
\end{lstlisting}
\paragraph{The bug.} \texttt{prev = 0} compares the first reading with a value that was never recorded. The program fails only when the first reading itself exceeds \texttt{T} while no real consecutive pair does.

\paragraph{Why the form matters.} \texttt{prev = 0} now reads as a heart rate of 0~bpm before the first reading, a value that makes no sense for a monitored patient. And any realistic recording exposes the bug, because a plausible heart rate is far larger than a plausible threshold.

\paragraph{Counter-example.} \texttt{has\_large\_change([72, 74], 15)}: the task requires \texttt{False}, the program returns \texttt{True}.

\paragraph{Valid Step 3 inputs.} Inputs must also be \emph{plausible in the setting}; students are told so, and an implausible Step~3 input is rejected with the message below and uses up an attempt.
\begin{lstlisting}[title={\small Validity check (server-side)}]
def valid(readings, T):
    if not isinstance(readings, list) or not all(type(r) is int for r in readings):
        return 'readings must be a list of whole numbers.'
    if type(T) is not int or T <= 0:
        return 'T must be a positive whole number.'
    if any(r < 30 or r > 220 for r in readings) or T > 60:
        return ('Your input must be a realistic recording: every heart rate '
                'between 30 and 220 bpm, and T at most 60 bpm.')
    return None
\end{lstlisting}
\paragraph{Expected effect.} Higher correctness than Direct; longer reading but a shorter search for an input.

\clearpage
\section{Problem P2: Warmest Night (Leh)}\label{sec:p2}
\subsection{Shared by all four variants}
The function signature, the reference solution, and the inputs used in Steps~1 and~2 are identical in all four variants of this problem. The Step~2 inputs are the same for both buggy programs, and neither program fails on any of them.

\paragraph{Signature.} \texttt{def warmest(temps):}

The Direct form names \texttt{warmest} as \texttt{largest}, \texttt{temps} as \texttt{nums}, in the signature, all programs and all inputs; nothing else changes.

\begin{lstlisting}[title={\small Reference solution (never shown to students)}]
def warmest(temps):
    best = temps[0]
    best_day = 1
    for i in range(len(temps)):
        if temps[i] > best:
            best = temps[i]
            best_day = i + 1
    return best_day
\end{lstlisting}
\paragraph{Step 2 tracing inputs.} The control-flow question (how many times each line runs) uses the first input; the data-flow question (variable values) uses the second. Reference outputs:

\smallskip\noindent\begin{tabular}{@{}l l@{}}\toprule
\textbf{Input} & \textbf{Output} \\\midrule
\texttt{warmest([2.5, 6.0, 4.0])} & \texttt{2} \\
\texttt{warmest([1.0, 3.5])} & \texttt{2} \\
\bottomrule\end{tabular}

\paragraph{Step 1 worked examples.} Revealed, with their correct outputs, only to a student who asks for help after a wrong Step~1 attempt; the student may not then reuse them as their own pairs:

\smallskip\noindent\begin{tabular}{@{}l l@{}}\toprule
\textbf{Input} & \textbf{Output} \\\midrule
\texttt{warmest([-2.0, -9.5, -11.0])} & \texttt{1} \\
\texttt{warmest([1.5, -6.0, 3.0])} & \texttt{3} \\
\bottomrule\end{tabular}

\subsection{Variant \texttt{P2\_Direct\_buggy1}}
\textbf{Form:} Direct. \quad \textbf{Bug:} Buggy1, context-independent.

\paragraph{Specification shown to the student (34 words).}
\begin{quote}\small Write \texttt{largest(nums)} that takes a non-empty list of numbers and returns the position of its largest value, counting positions from 1. If the largest value occurs more than once, return the smallest such position.\end{quote}
\begin{lstlisting}[title={\small Program shown in Steps 2 and 3}]
def largest(nums):
    best = nums[0]
    best_day = 1
    for i in range(len(nums)):
        if nums[i] >= best:      # <- bug
            best = nums[i]
            best_day = i + 1
    return best_day
\end{lstlisting}
\paragraph{The bug.} Because it uses \texttt{>=}, it moves to every later value equal to the best so far, so when the largest value occurs more than once it returns the last position instead of the first. A counter-example needs a tie at the maximum.

\paragraph{Why the form matters.} The tie rule (smallest position) is stated explicitly, exactly as in the contextualized form.

\paragraph{Counter-example.} \texttt{warmest([-9.8, -12.0, -9.8])}: the task requires \texttt{1}, the program returns \texttt{3}.

\paragraph{Valid Step 3 inputs.} Inputs must satisfy the specification's own preconditions.
\begin{lstlisting}[title={\small Validity check (server-side)}]
def valid(temps):
    if not isinstance(temps, list) or len(temps) == 0:
        return 'temps must be a non-empty list of numbers.'
    if not all(type(t) in (int, float) for t in temps):
        return 'temps must be a non-empty list of numbers.'
    return None
\end{lstlisting}
\paragraph{Expected effect.} Baseline for this bug.

\subsection{Variant \texttt{P2\_Contextualized\_buggy1}}
\textbf{Form:} Contextualized. \quad \textbf{Bug:} Buggy1, context-independent.

\paragraph{Specification shown to the student (95 words).}
\begin{quote}\small A weather station in Leh, one of the coldest inhabited places in India, records the minimum temperature of each night in January, in degrees Celsius to one decimal place. The list has at least 1 entry, starts with the night of 1 January and has one entry per night, up to the most recent night recorded. The district administration wants to know which night was the warmest. Write \texttt{warmest(temps)} that takes the list and returns the date in January of the warmest night. If more than one night was equally warm, return the earliest date.\end{quote}
\begin{lstlisting}[title={\small Program shown in Steps 2 and 3}]
def warmest(temps):
    best = temps[0]
    best_day = 1
    for i in range(len(temps)):
        if temps[i] >= best:      # <- bug
            best = temps[i]
            best_day = i + 1
    return best_day
\end{lstlisting}
\paragraph{The bug.} Because it uses \texttt{>=}, it moves to every later value equal to the best so far, so when the largest value occurs more than once it returns the last position instead of the first. A counter-example needs a tie at the maximum.

\paragraph{Why the form matters.} The tie rule (earliest date) is stated explicitly, as in the direct form. The setting may still matter: temperatures to one decimal place make exact ties seem less likely to a student imagining realistic data.

\paragraph{Counter-example.} \texttt{warmest([-9.8, -12.0, -9.8])}: the task requires \texttt{1}, the program returns \texttt{3}.

\paragraph{Valid Step 3 inputs.} Inputs must also be \emph{plausible in the setting}; students are told so, and an implausible Step~3 input is rejected with the message below and uses up an attempt.
\begin{lstlisting}[title={\small Validity check (server-side)}]
def valid(temps):
    if not isinstance(temps, list) or len(temps) == 0:
        return 'temps must be a non-empty list of numbers.'
    if not all(type(t) in (int, float) for t in temps):
        return 'temps must be a non-empty list of numbers.'
    if (len(temps) > 31 or any(t < -40 or t > 10 for t in temps)
            or any(round(t, 1) != t for t in temps)):
        return ('Your input must be a realistic record: at most 31 nights, '
                'each temperature between -40.0 and 10.0 degrees Celsius, '
                'to one decimal place.')
    return None
\end{lstlisting}
\paragraph{Expected effect.} No gain over Direct; possibly lower for less fluent readers. Longer time (more to read).

\subsection{Variant \texttt{P2\_Direct\_buggy2}}
\textbf{Form:} Direct. \quad \textbf{Bug:} Buggy2, context-sensitive.

\paragraph{Specification shown to the student (34 words).}
\begin{quote}\small Write \texttt{largest(nums)} that takes a non-empty list of numbers and returns the position of its largest value, counting positions from 1. If the largest value occurs more than once, return the smallest such position.\end{quote}
\begin{lstlisting}[title={\small Program shown in Steps 2 and 3}]
def largest(nums):
    best = 0                     # <- bug
    best_day = 1
    for i in range(len(nums)):
        if nums[i] > best:
            best = nums[i]
            best_day = i + 1
    return best_day
\end{lstlisting}
\paragraph{The bug.} \texttt{best = 0} means a value replaces it only if it is above zero. When every value is at or below zero, \texttt{best} never changes and the program returns 1.

\paragraph{Why the form matters.} With only ``numbers'' to go on, 0 is a natural-looking starting value, and the program is correct on the positive lists a student is likely to try: on \texttt{[3.0, 7.5, 5.0]} it correctly returns 2.

\paragraph{Counter-example.} \texttt{warmest([-14.2, -9.8])}: the task requires \texttt{2}, the program returns \texttt{1}.

\paragraph{Valid Step 3 inputs.} Inputs must satisfy the specification's own preconditions.
\begin{lstlisting}[title={\small Validity check (server-side)}]
def valid(temps):
    if not isinstance(temps, list) or len(temps) == 0:
        return 'temps must be a non-empty list of numbers.'
    if not all(type(t) in (int, float) for t in temps):
        return 'temps must be a non-empty list of numbers.'
    return None
\end{lstlisting}
\paragraph{Expected effect.} Baseline for this bug.

\subsection{Variant \texttt{P2\_Contextualized\_buggy2}}
\textbf{Form:} Contextualized. \quad \textbf{Bug:} Buggy2, context-sensitive.

\paragraph{Specification shown to the student (95 words).}
\begin{quote}\small A weather station in Leh, one of the coldest inhabited places in India, records the minimum temperature of each night in January, in degrees Celsius to one decimal place. The list has at least 1 entry, starts with the night of 1 January and has one entry per night, up to the most recent night recorded. The district administration wants to know which night was the warmest. Write \texttt{warmest(temps)} that takes the list and returns the date in January of the warmest night. If more than one night was equally warm, return the earliest date.\end{quote}
\begin{lstlisting}[title={\small Program shown in Steps 2 and 3}]
def warmest(temps):
    best = 0                      # <- bug
    best_day = 1
    for i in range(len(temps)):
        if temps[i] > best:
            best = temps[i]
            best_day = i + 1
    return best_day
\end{lstlisting}
\paragraph{The bug.} \texttt{best = 0} means a value replaces it only if it is above zero. When every value is at or below zero, \texttt{best} never changes and the program returns 1.

\paragraph{Why the form matters.} Every realistic January record from Leh is below zero, so any record whose warmest night is not 1~January exposes the bug. And \texttt{best = 0} now says the warmest night so far was 0\,\textdegree C, warmer than any night in the record, a value that can never be beaten.

\paragraph{Counter-example.} \texttt{warmest([-14.2, -9.8])}: the task requires \texttt{2}, the program returns \texttt{1}.

\paragraph{Valid Step 3 inputs.} Inputs must also be \emph{plausible in the setting}; students are told so, and an implausible Step~3 input is rejected with the message below and uses up an attempt.
\begin{lstlisting}[title={\small Validity check (server-side)}]
def valid(temps):
    if not isinstance(temps, list) or len(temps) == 0:
        return 'temps must be a non-empty list of numbers.'
    if not all(type(t) in (int, float) for t in temps):
        return 'temps must be a non-empty list of numbers.'
    if (len(temps) > 31 or any(t < -40 or t > 10 for t in temps)
            or any(round(t, 1) != t for t in temps)):
        return ('Your input must be a realistic record: at most 31 nights, '
                'each temperature between -40.0 and 10.0 degrees Celsius, '
                'to one decimal place.')
    return None
\end{lstlisting}
\paragraph{Expected effect.} Higher correctness than Direct; longer reading but a shorter search for an input.

\clearpage
\section{Database fields of a problem}\label{sec:db}
Each variant is stored as one row of the \texttt{tasks} table. Table~\ref{tab:fields} explains every field and says whether, and when, a student sees it; Table~\ref{tab:example} shows the values stored for one variant. Fields marked \emph{Never} are not sent to the browser at all.

{\small
\begin{longtable}{@{}>{\raggedright\arraybackslash}p{3.6cm} >{\raggedright\arraybackslash}p{1.7cm} >{\raggedright\arraybackslash}p{2.0cm} >{\raggedright\arraybackslash}p{7.3cm}@{}}
\caption{Fields of one problem (\texttt{tasks} table).}\label{tab:fields}\\
\toprule\textbf{Field} & \textbf{Type} & \textbf{Seen by student} & \textbf{Meaning} \\
\endfirsthead
\toprule\textbf{Field} & \textbf{Type} & \textbf{Seen by student} & \textbf{Meaning} \\
\endhead
\bottomrule\endfoot
\midrule\multicolumn{4}{@{}l}{\textit{Identity}} \\\addlinespace[2pt]
\texttt{id} & integer & Yes (opaque) & Database key, assigned on insert. \\
\texttt{title} & text & Never & Unique name; for study tasks equal to the variant key. Re-seeding matches existing rows by title, so it never duplicates. Study sessions send the browser a neutral ``Problem 1'' or ``Problem 2'' instead. \\
\texttt{language} & text & No & Programming language; always \texttt{python}. \\
\midrule\multicolumn{4}{@{}l}{\textit{Study design}} \\\addlinespace[2pt]
\texttt{variant\_\allowbreak{}key} & text & Never & Which of the eight variants this row is, e.g.\ \texttt{P2\_Contextualized\_buggy2}. Used by the assignment; null for practice tasks. \\
\texttt{problem\_\allowbreak{}key} & text & Never & \texttt{P1} or \texttt{P2}. \\
\texttt{form} & text & Yes (Step 3) & \texttt{direct} or \texttt{contextualized}. Sent to the browser only to decide whether to show the ``input must be realistic'' notice. \\
\texttt{bug\_\allowbreak{}kind} & text & Never & \texttt{buggy1} (context-independent) or \texttt{buggy2} (context-sensitive). Analysis only. \\
\midrule\multicolumn{4}{@{}l}{\textit{Specification}} \\\addlinespace[2pt]
\texttt{description} & text & Steps 1, 3 & The English specification (the form's text). \\
\texttt{description\_\allowbreak{}mixed} & JSON & Steps 1, 3 & Code-mixed versions keyed by native-language code (\texttt{te}, \texttt{ta}, \texttt{hi}, \texttt{gu}, \texttt{de}). A student receives only their own language; with no entry, no toggle is shown. \\
\texttt{function\_\allowbreak{}signature} & text & Steps 1, 3 & First line of the function. The function name is fixed in the answer boxes, so students type only the arguments. \\
\midrule\multicolumn{4}{@{}l}{\textit{Programs}} \\\addlinespace[2pt]
\texttt{correct\_\allowbreak{}code} & text & Never & Reference solution. Step~1 pairs and Step~3 required outputs are checked against it. \\
\texttt{buggy\_\allowbreak{}code} & text & Steps 2, 3 & The faulty program for this variant. \\
\texttt{buggy\_\allowbreak{}line\_\allowbreak{}numbers} & JSON list & Never & Line(s) carrying the bug. Lets the analysis ask whether the student traced the faulty line correctly. \\
\midrule\multicolumn{4}{@{}l}{\textit{Step 1}} \\\addlinespace[2pt]
\texttt{io\_\allowbreak{}example\_\allowbreak{}inputs} & JSON list & After a wrong attempt & Two calls whose correct outputs are revealed to a student who asks for help; then barred as the student's own pairs. \\
\midrule\multicolumn{4}{@{}l}{\textit{Step 2}} \\\addlinespace[2pt]
\texttt{trace\_\allowbreak{}sample\_\allowbreak{}inputs} & JSON list & Step 2 & Call(s) for the control-flow question (how many times each line runs). Never a counter-example. \\
\texttt{trace\_\allowbreak{}data\_\allowbreak{}flow\_\allowbreak{}inputs} & JSON list & Step 2 & Different call(s) for the data-flow question (variable values, with the execution order given). \\
\texttt{trace\_\allowbreak{}prefill\_\allowbreak{}steps} & integer & Step 2 & Rows filled in as a worked example (1). \\
\texttt{trace\_\allowbreak{}omit\_\allowbreak{}unchanged\_\allowbreak{}vars} & boolean & No & Whether to drop never-changing variables from the value table (off). \\
\texttt{mutation\_\allowbreak{}line\_\allowbreak{}number, mutation\_\allowbreak{}new\_\allowbreak{}line\_\allowbreak{}text, mutation\_\allowbreak{}prompt} & int, text, text & --- & Optional reflection question from the practice tasks. Empty for all study variants. \\
\midrule\multicolumn{4}{@{}l}{\textit{Step 3}} \\\addlinespace[2pt]
\texttt{input\_\allowbreak{}validator} & text (Python) & Its messages & Defines \texttt{valid(...)} with the task's parameters; returns \texttt{None} for a valid input or the message shown to the student. The contextualized version adds plausibility ranges. \\
\texttt{step3\_\allowbreak{}max\_\allowbreak{}attempts} & integer & Yes & Attempt limit (3). Set for study tasks only; it also switches Step~3 to the study format (input, required output, program output). \\
\end{longtable}
}
\begin{table}[H]\centering\small
\caption{The stored values of \texttt{P2\_Contextualized\_buggy2} (long text fields omitted).}
\label{tab:example}
\begin{tabularx}{\textwidth}{@{}l X@{}}\toprule
\textbf{Field} & \textbf{Value} \\\midrule
\texttt{title / variant\_key} & \texttt{P2\_Contextualized\_buggy2} \\
\texttt{problem\_key} & \texttt{P2} \\
\texttt{form} & \texttt{contextualized} \\
\texttt{bug\_kind} & \texttt{buggy2} \\
\texttt{function\_signature} & \texttt{def warmest(temps):} \\
\texttt{buggy\_line\_numbers} & \texttt{[2]} \\
\texttt{trace\_sample\_inputs} & \texttt{['warmest([2.5, 6.0, 4.0])']} \\
\texttt{trace\_data\_flow\_inputs} & \texttt{['warmest([1.0, 3.5])']} \\
\texttt{io\_example\_inputs} & \texttt{['warmest([-2.0, -9.5, -11.0])', 'warmest([1.5, -6.0, 3.0])']} \\
\texttt{description\_mixed} & \texttt{\{te: ..., ta: ..., hi: ..., gu: ..., de: ...\}} \\
\texttt{step3\_max\_attempts} & \texttt{3} \\
\texttt{mutation\_*} & \texttt{null} \\
\bottomrule\end{tabularx}\end{table}

\subsection{Related tables}
A problem row is the fixed material. What students do is recorded in four
other tables:

\begin{description}[style=nextline, leftmargin=1.5em]
  \item[\texttt{participants}] One row per student: institute, student ID,
    the questionnaire answers (\texttt{native\_\allowbreak{}language},
    \texttt{schooling\_\allowbreak{}medium}, \texttt{english\_\allowbreak{}proficiency},
    \texttt{mother\_\allowbreak{}language\_\allowbreak{}proficiency},
    \texttt{programming\_\allowbreak{}proficiency},
    \texttt{programming\_\allowbreak{}courses}, \texttt{age}, \texttt{gender}; see
    Section~\ref{sec:questionnaire}), arrival \texttt{slot} at the institute, the
    \texttt{condition\_row} it was dealt, and \texttt{assigned\_task\_ids}
    (the two variants, in presentation order). Institute and student ID are
    unique together, so a returning student resumes rather than being
    re-assigned.
  \item[\texttt{study\_sessions}] One row per student $\times$ problem:
    \texttt{participant\_id}, \texttt{position} (1 or 2), \texttt{task\_id},
    and whether Steps 1--3 are completed.
  \item[\texttt{step\_events}] The timing and answer log: one row each time a
    step is shown and each time it is submitted, with a server timestamp,
    attempt number, correctness, and the full submission in
    \texttt{payload}. Step~1 payloads also record, for analysis only,
    whether each pair would already have refuted the program; Step~3
    payloads record the real outputs and each success condition separately.
  \item[\texttt{spec\_view\_events}] One row for each step's first view of the
    specification (always English) and one per switch between English and
    code-mixed, with a timestamp.
\end{description}
\clearpage
\appendix
\section{Code-mixed specifications}\label{app:mixed}
\textbf{Draft.} These versions must be checked by a native speaker of each language against the English text before the pilot; an ambiguity introduced in translation would produce disputed counter-examples. Technical terms, identifiers and code stay in English; the surrounding prose is in the student's language, written in its own script (German: German prose with English technical terms).

\subsection{P1, Direct form}
\paragraph{Telugu.}
{\tefont \texttt{has\_large\_change(readings, T)} అనే function రాయండి. ఇది integers ఉన్న \texttt{readings} అనే list ని, ఒక positive integer \texttt{T} ని తీసుకుంటుంది. \texttt{readings} లో ఏదైనా ఒక consecutive elements pair మధ్య absolute difference \texttt{T} కంటే greater అయితే \texttt{True} return చేస్తుంది, లేకపోతే \texttt{False} return చేస్తుంది. List empty గా ఉండొచ్చు, లేదా అందులో ఒకే ఒక element ఉండొచ్చు.}

\paragraph{Tamil.}
{\tafont \texttt{has\_large\_change(readings, T)} என்ற function-ஐ எழுதுங்கள். இது integers கொண்ட \texttt{readings} என்ற list-ஐயும், ஒரு positive integer \texttt{T}-ஐயும் எடுத்துக்கொள்ளும். \texttt{readings}-இல் ஏதாவது ஒரு consecutive elements pair-க்கு இடையிலான absolute difference \texttt{T}-ஐ விட greater ஆக இருந்தால் \texttt{True} return செய்யும், இல்லையென்றால் \texttt{False} return செய்யும். List empty ஆக இருக்கலாம், அல்லது அதில் ஒரே ஒரு element மட்டும் இருக்கலாம்.}

\paragraph{Hindi.}
{\hifont एक function \texttt{has\_large\_change(readings, T)} लिखिए जो integers की एक list \texttt{readings} और एक positive integer \texttt{T} लेता है। अगर \texttt{readings} के किसी भी consecutive elements के pair का absolute difference \texttt{T} से greater है, तो यह \texttt{True} return करता है, वरना \texttt{False}। List empty हो सकती है या उसमें सिर्फ़ एक element हो सकता है।}

\paragraph{Gujarati.}
{\gufont \texttt{has\_large\_change(readings, T)} નામનું function લખો, જે integers ની એક list \texttt{readings} અને એક positive integer \texttt{T} લે છે. જો \texttt{readings} ના કોઈ પણ consecutive elements ની pair વચ્ચેનો absolute difference \texttt{T} કરતાં greater હોય, તો તે \texttt{True} return કરે છે, નહીં તો \texttt{False}. List empty હોઈ શકે છે અથવા તેમાં ફક્ત એક જ element હોઈ શકે છે.}

\paragraph{German.}
{Schreiben Sie eine Funktion \texttt{has\_large\_change(readings, T)}, die eine Liste \texttt{readings} von Integer-Werten und einen positiven Integer-Wert \texttt{T} entgegennimmt. Wenn die absolute Differenz zwischen irgendeinem Paar aufeinanderfolgender Elemente von \texttt{readings} größer als \texttt{T} ist, gibt sie \texttt{True} zurück, andernfalls \texttt{False}. Die Liste kann leer sein oder nur ein Element enthalten.}

\subsection{P1, Contextualized form}
\paragraph{Telugu.}
{\tefont ఒక hospital లోని patient monitor, patient యొక్క heart rate ని beats per minute (bpm) లో, ప్రతి second కి ఒకసారి record చేస్తుంది. Healthy heart rate ఒక second నుండి తర్వాతి second కి కొంచెం మాత్రమే మారుతుంది. ఏవైనా రెండు consecutive readings మధ్య heart rate ఎప్పుడైనా \texttt{T} bpm కంటే ఎక్కువ మారితే, nurse ఆ recording ని review చేయాలి అని hospital clinical guidelines చెబుతున్నాయి: sudden drop అంటే heart fail అవుతోందని అర్థం కావచ్చు, sudden rise అంటే patient distress లో ఉన్నారని అర్థం కావచ్చు. Recording empty గా ఉండొచ్చు లేదా అందులో ఒకే ఒక reading ఉండొచ్చు, ఉదాహరణకు monitor ని ఇప్పుడే attach చేసినప్పుడు. Readings list (whole numbers) ని, threshold \texttt{T} (ఒక positive whole number) ని తీసుకుని, recording కి review అవసరమైతే \texttt{True}, లేకపోతే \texttt{False} return చేసే \texttt{has\_large\_change(readings, T)} అనే function రాయండి.}

\paragraph{Tamil.}
{\tafont ஒரு hospital-இன் patient monitor, ஒரு patient-இன் heart rate-ஐ beats per minute (bpm)-இல், ஒவ்வொரு second-க்கும் ஒரு முறை record செய்கிறது. Healthy heart rate ஒரு second-இலிருந்து அடுத்த second-க்கு கொஞ்சம் மட்டுமே மாறும். ஏதாவது இரண்டு consecutive readings-க்கு இடையில் heart rate எப்போதாவது \texttt{T} bpm-ஐ விட அதிகமாக மாறினால், nurse அந்த recording-ஐ review செய்ய வேண்டும் என்று hospital-இன் clinical guidelines சொல்கின்றன: திடீர் drop என்றால் heart fail ஆகிறது என்று அர்த்தமாக இருக்கலாம், திடீர் rise என்றால் patient distress-இல் இருக்கிறார் என்று அர்த்தமாக இருக்கலாம். Recording empty ஆக இருக்கலாம் அல்லது அதில் ஒரே ஒரு reading மட்டும் இருக்கலாம், உதாரணமாக monitor இப்போதுதான் attach செய்யப்பட்டிருந்தால். Readings list (whole numbers) மற்றும் threshold \texttt{T} (ஒரு positive whole number) ஆகியவற்றை எடுத்துக்கொண்டு, recording-க்கு review தேவைப்பட்டால் \texttt{True}, இல்லையென்றால் \texttt{False} return செய்யும் \texttt{has\_large\_change(readings, T)} என்ற function-ஐ எழுதுங்கள்.}

\paragraph{Hindi.}
{\hifont एक hospital का patient monitor हर second एक बार patient का heart rate, beats per minute (bpm) में, record करता है। एक healthy heart rate एक second से अगले second तक बस थोड़ा-सा ही बदलता है। Hospital की clinical guidelines कहती हैं कि अगर किन्हीं भी दो consecutive readings के बीच heart rate कभी भी \texttt{T} bpm से ज़्यादा बदलता है, तो nurse को उस recording को review करना होगा: अचानक drop का मतलब हो सकता है कि heart fail हो रहा है, और अचानक rise का मतलब हो सकता है कि patient distress में है। Recording empty हो सकती है या उसमें सिर्फ़ एक reading हो सकती है, जैसे अगर monitor अभी-अभी लगाया गया हो। एक function \texttt{has\_large\_change(readings, T)} लिखिए जो readings की list (whole numbers) और threshold \texttt{T} (एक positive whole number) लेता है, और अगर recording को review की ज़रूरत है तो \texttt{True} return करता है, वरना \texttt{False}।}

\paragraph{Gujarati.}
{\gufont એક hospital નું patient monitor દર second એક વાર patient નો heart rate, beats per minute (bpm) માં, record કરે છે. Healthy heart rate એક second થી બીજી second સુધી ફક્ત થોડો જ બદલાય છે. Hospital ની clinical guidelines કહે છે કે જો કોઈ પણ બે consecutive readings વચ્ચે heart rate ક્યારેય \texttt{T} bpm કરતાં વધારે બદલાય, તો nurse એ તે recording ને review કરવું જ પડે: અચાનક drop નો અર્થ હોઈ શકે કે heart fail થઈ રહ્યું છે, અને અચાનક rise નો અર્થ હોઈ શકે કે patient distress માં છે. Recording empty હોઈ શકે છે અથવા તેમાં ફક્ત એક જ reading હોઈ શકે છે, ઉદાહરણ તરીકે જો monitor હમણાં જ attach કર્યું હોય. \texttt{has\_large\_change(readings, T)} નામનું function લખો, જે readings ની list (whole numbers) અને threshold \texttt{T} (એક positive whole number) લે છે, અને જો recording ને review ની જરૂર હોય તો \texttt{True} return કરે છે, નહીં તો \texttt{False}.}

\paragraph{German.}
{Der Patientenmonitor eines Krankenhauses erfasst einmal pro Sekunde die Herzfrequenz eines Patienten, gemessen in beats per minute (bpm). Eine gesunde Herzfrequenz ändert sich von einer Sekunde zur nächsten nur geringfügig. Die klinischen Richtlinien des Krankenhauses besagen, dass eine Pflegekraft eine Aufzeichnung überprüfen muss, wenn sich die Herzfrequenz zwischen zwei aufeinanderfolgenden Messwerten jemals um mehr als \texttt{T} bpm ändert: Ein plötzlicher Abfall kann bedeuten, dass das Herz versagt, und ein plötzlicher Anstieg kann bedeuten, dass der Patient in Not ist. Eine Aufzeichnung kann leer sein oder nur einen einzelnen Messwert enthalten, zum Beispiel wenn der Monitor gerade erst angebracht wurde. Schreiben Sie eine Funktion \texttt{has\_large\_change(readings, T)}, die die Liste der Messwerte (ganze Zahlen) und den Schwellenwert \texttt{T} (eine positive ganze Zahl) entgegennimmt und \texttt{True} zurückgibt, wenn die Aufzeichnung überprüft werden muss, und andernfalls \texttt{False}.}

\subsection{P2, Direct form}
\paragraph{Telugu.}
{\tefont Numbers ఉన్న ఒక non-empty list ని తీసుకుని, దానిలోని largest value యొక్క position ని return చేసే \texttt{largest(nums)} రాయండి; positions ని 1 నుండి లెక్కించాలి. Largest value ఒకటి కంటే ఎక్కువ సార్లు వస్తే, వాటిలో smallest position ని return చేయండి.}

\paragraph{Tamil.}
{\tafont Numbers கொண்ட ஒரு non-empty list-ஐ எடுத்துக்கொண்டு, அதன் largest value-இன் position-ஐ return செய்யும் \texttt{largest(nums)}-ஐ எழுதுங்கள்; positions-ஐ 1-இலிருந்து எண்ண வேண்டும். Largest value ஒன்றுக்கு மேற்பட்ட முறை வந்தால், அவற்றில் smallest position-ஐ return செய்யுங்கள்.}

\paragraph{Hindi.}
{\hifont \texttt{largest(nums)} लिखिए जो numbers की एक non-empty list लेता है और उसकी largest value की position return करता है, positions को 1 से गिनते हुए। अगर largest value एक से ज़्यादा बार आती है, तो ऐसी सबसे छोटी (smallest) position return कीजिए।}

\paragraph{Gujarati.}
{\gufont \texttt{largest(nums)} લખો, જે numbers ની એક non-empty list લે છે અને તેની largest value ની position return કરે છે; positions ને 1 થી ગણવાની છે. જો largest value એક કરતાં વધુ વાર આવે, તો એમાંથી સૌથી નાની (smallest) position return કરો.}

\paragraph{German.}
{Schreiben Sie eine Funktion \texttt{largest(nums)}, die eine nicht-leere Liste \texttt{nums} von Zahlen entgegennimmt und die Position ihres größten Werts zurückgibt, wobei die Positionen ab 1 gezählt werden. Kommt der größte Wert mehrfach vor, geben Sie die kleinste dieser Positionen zurück.}

\subsection{P2, Contextualized form}
\paragraph{Telugu.}
{\tefont India లో మనుషులు నివసించే అత్యంత చల్లని ప్రదేశాలలో ఒకటైన Leh లోని ఒక weather station, January లో ప్రతి రాత్రి యొక్క lowest temperature ని degrees Celsius లో, ఒక decimal place వరకు record చేస్తుంది. List లో కనీసం 1 entry ఉంటుంది; అది 1 January రాత్రితో మొదలవుతుంది, ప్రతి రాత్రికి ఒక entry ఉంటుంది, చివరిగా record చేసిన రాత్రి వరకు. ఏ రాత్రి అత్యంత warm గా ఉందో district administration తెలుసుకోవాలనుకుంటోంది. List ని తీసుకుని, అత్యంత warm రాత్రి యొక్క January date ని return చేసే \texttt{warmest(temps)} రాయండి. ఒకటి కంటే ఎక్కువ రాత్రులు సమానంగా warm గా ఉంటే, earliest date ని return చేయండి.}

\paragraph{Tamil.}
{\tafont India-வில் மக்கள் வாழும் மிகக் குளிரான இடங்களில் ஒன்றான Leh-இல் உள்ள ஒரு weather station, January மாதத்தில் ஒவ்வொரு இரவின் lowest temperature-ஐ degrees Celsius-இல், ஒரு decimal place வரை record செய்கிறது. List-இல் குறைந்தது 1 entry இருக்கும்; அது 1 January இரவில் தொடங்கி, ஒவ்வொரு இரவுக்கும் ஒரு entry என்று, கடைசியாக record செய்யப்பட்ட இரவு வரை இருக்கும். எந்த இரவு மிகவும் warm ஆக இருந்தது என்று district administration தெரிந்துகொள்ள விரும்புகிறது. List-ஐ எடுத்துக்கொண்டு, மிகவும் warm ஆன இரவின் January date-ஐ return செய்யும் \texttt{warmest(temps)}-ஐ எழுதுங்கள். ஒன்றுக்கு மேற்பட்ட இரவுகள் சமமாக warm ஆக இருந்தால், earliest date-ஐ return செய்யுங்கள்.}

\paragraph{Hindi.}
{\hifont Leh में, जो India की सबसे ठंडी आबादी वाली जगहों में से एक है, एक weather station January में हर रात का lowest temperature, degrees Celsius में, एक decimal place तक record करता है। List में कम से कम 1 entry है, यह 1 January की रात से शुरू होती है और हर रात की एक entry है, सबसे recent record की गई रात तक। District administration जानना चाहता है कि कौन-सी रात सबसे warm थी। \texttt{warmest(temps)} लिखिए जो list लेता है और सबसे warm रात की January की date return करता है। अगर एक से ज़्यादा रातें बराबर warm थीं, तो सबसे पहली (earliest) date return कीजिए।}

\paragraph{Gujarati.}
{\gufont India ની સૌથી ઠંડી વસવાટવાળી જગ્યાઓમાંની એક, Leh માં એક weather station January ની દરેક રાતનું lowest temperature, degrees Celsius માં, એક decimal place સુધી record કરે છે. List માં ઓછામાં ઓછી 1 entry છે; તે 1 January ની રાતથી શરૂ થાય છે અને દરેક રાત માટે એક entry છે, છેલ્લે record થયેલી રાત સુધી. District administration જાણવા માંગે છે કે કઈ રાત સૌથી warm હતી. \texttt{warmest(temps)} લખો, જે list લે છે અને સૌથી warm રાતની January ની date return કરે છે. જો એક કરતાં વધુ રાતો સરખી warm હોય, તો સૌથી પહેલી (earliest) date return કરો.}

\paragraph{German.}
{Eine Wetterstation in Kühnhaide / Sachsen, einem der kältesten bewohnten Orte Deutschlands, zeichnet die Tiefsttemperatur jeder Januarnacht in Grad Celsius mit einer Dezimalstelle auf. Die Liste enthält mindestens 1 Eintrag, beginnt mit der Nacht vom 1. Januar und enthält einen Eintrag pro Nacht, bis zur zuletzt aufgezeichneten Nacht. Die Bezirksverwaltung möchte wissen, welche Nacht die wärmste war. Schreiben Sie eine Funktion \texttt{warmest(temps)}, die die Liste entgegennimmt und das Datum im Januar der wärmsten Nacht zurückgibt. Wenn mehrere Nächte gleich warm waren, geben Sie das früheste Datum zurück.}


\end{document}
